% !TeX root = ../深度学习与大语言模型技术原理.tex
% 由主文档通过 \input 载入；请编译主文档。
\setcounter{chapter}{7} % 保留原书章号。
\chapter{训练稳定性、精度与显存}
\label{ch:stability}

\section{参数之外还有哪些显存}

训练峰值显存近似由参数、梯度、优化器状态、前向激活、临时缓冲和通信缓冲组成。以一种混合精度 Adam 实现为例，每参数存2字节低精度权重、2字节梯度、4字节主权重和两个4字节矩估计，静态状态约为 $16P$ 字节。若实现没有独立主权重、使用32位梯度或分片状态，系数就不同；这不是所有框架的固定账单。

激活内存与批量、序列长度、层数及中间宽度有关。保存注意力矩阵的实现还包含二次长度项；优化的注意力内核会改变该项。参数量很少变化时，增加上下文长度仍可能导致显存不足，所以必须记录峰值及其发生阶段，区分前向、反向、优化器更新和评测生成。

\begin{example}
一个模型有 $P=10^9$ 个参数，采用上述每参数16字节状态估算，四卡均匀分片后每卡静态状态是否就是训练所需全部显存？
\end{example}
\begin{solve}
总静态状态约为16 GB，理想全分片后每卡约4 GB，但还必须加激活、当前聚合的完整参数块、通信及算子临时空间。分片只能减少被分片的状态，不能据此认定4 GB显存即可训练。
\end{solve}

\section{混合精度与损失缩放}

混合精度（mixed precision）在存储、矩阵乘法、归约和优化器更新中采用不同数据类型。FP16 的指数位较少，BF16 保留与 FP32 相同数量的指数位但尾数更短，因此 BF16 的动态范围更宽而相对精度更低。它们均可能出现舍入误差；某些归约、归一化或 softmax 运算需要更高精度累积。FP8 等格式进一步降低位数，但尺度管理、硬件支持和算子覆盖更为关键。

若低精度梯度易下溢，可以把损失乘尺度 $s>0$ 后反向传播，再把得到的梯度除以 $s$。这在实数算术中等价，但有限精度下能让较小梯度落入可表示区间。动态损失缩放在检测到非有限梯度时跳过更新并降低尺度；必须先反缩放再按真实梯度尺度裁剪。BF16 常常不依赖损失缩放，但不能因此忽视溢出和非有限值检查。

定位损失尖峰时，应沿输入、激活、logits、损失、梯度和参数更新逐级检查有限性。数据中的异常长度、错误标签、过大的学习率、归一化失稳与内核错误都可能产生类似现象。直接减小学习率有时能缓解症状，却不能替代定位根因。

\section{梯度累积与有效词元归一化}

梯度累积（gradient accumulation）把一个有效批量拆成多个微批量（microbatch），多次反向后才更新参数。若第 $j$ 个微批量损失总和为 $S_j$、有效词元数为 $N_j$，希望优化的是
\begin{equation}
\mathcal L_{\mathrm{batch}}=\frac{\sum_jS_j}{\sum_jN_j}.
\label{eq:token-normalization}
\end{equation}
先对每个微批量求平均再等权累加，仅在所有 $N_j$ 相等时与此一致。变长序列、回答区域掩码和多设备有效词元不等时，这个区别尤其重要。

在固定参数下、损失可按样本分解、没有依赖批量统计的操作且随机性处理相容时，累积梯度等于同一大批量的梯度。浮点求和顺序、dropout 和批归一化等因素会破坏逐位相同。优化器和学习率调度器应按实际更新次数推进，不能每个微批量都更新一次。

全局范数裁剪把 $g$ 替换为 $g\min(1,c/\|g\|_2)$，$g=0$ 时保持不变。分片训练需要计算跨设备的全局范数，不能把每张卡独立裁剪误当成全局裁剪。裁剪限制单次梯度幅度，但不会修复错误目标或消除所有训练不稳定。

\section{激活检查点与计算换空间}

激活检查点（activation checkpointing）只保存部分前向边界激活，反向时重算内部激活。其输出梯度在相同函数与随机状态下应与完整保存相符；存在随机算子时，需要恢复相应随机数状态。含有不可重复副作用的函数不适合直接重计算。

为理解代价，设 $L$ 层每层激活大小约为 $A$，每隔 $s$ 层保存检查点。忽略其他状态时，边界保存约 $LA/s$，单段重算时额外保存约 $sA$，总量约 $A(L/s+s)$，在 $s\approx\sqrt L$ 附近较小。真实模型层大小不等、重算策略多样，这只是均匀链式网络的示意估算。节省空间需要额外前向计算，应测量每训练词元耗时是否仍满足预算。

